% Copyright 2026 hdl-simbench authors.
% SPDX-License-Identifier: Apache-2.0
%
% The hdl-simbench report. Every number in it comes from generated files:
% tables.tex and provenance.tex (//tools:report_tables, from each
% campaign's summary.csv) and version.tex (the release tag, or
% "unreleased"). Do not type numbers into this file.
\documentclass[11pt]{article}
\usepackage[margin=2cm]{geometry}
\usepackage{graphicx}
% lscape, not pdflscape, which the hermetic TinyTeX lacks: the landscape
% pages are set sideways on portrait pages.
\usepackage{lscape}

\input{version}

\title{hdl-simbench: open-source HDL simulators on the same designs}
\author{HDL Factory}
\date{\ReportVersion}

\begin{document}
\maketitle

\section{What this measures}

hdl-simbench runs the same designs under several open-source simulators,
built the same way, on the same machine, and times them. SystemVerilog
designs run under Verilator, Icarus Verilog and xezim; VHDL designs run
under GHDL and NVC. Simulators are compared only within a language.

Commercial simulators are not in the study: their licenses forbid
publishing benchmark results.

The VHDL case runs CoreMark\textsuperscript{\textregistered}, a
trademark of EEMBC, on the NEORV32 RISC-V processor as a workload for the
simulators. It measures the simulators, not the simulated processor, and
its numbers are not CoreMark scores.

\section{How the simulators are built}

Every simulator and every tool that builds one comes from a pinned
version through Bazel. A simulator binary may take only glibc from the
machine it runs on, except for host libraries a campaign explicitly
allowed; every run's command record names each allowed library with its
path and checksum, and the campaign notes say why.
C and C++ are compiled by clang from a pinned LLVM release against a
pinned Debian sysroot (glibc 2.31), optimized and without debug
information. Each timed run records every simulator command line, and each
simulator binary's dynamic dependencies are checked on every run; a binary
that needs anything else from the host fails the run. The
repository's \texttt{README.md} and \texttt{bench/vhdl/README.md} list
every compiler and simulator flag.

\section{How a run counts}

Timed runs pin the simulation to one physical core, with its
hyperthread sibling idle, on a machine whose CI runners are paused for the
campaign. Around each run the bench records the machine's load, the busy
share of the pinned core, of its sibling and of the other cores, and CPU
steal.

These rules were fixed before any result was seen, and apply unchanged:
\begin{itemize}
\item A run is dropped if its pinned core's sibling was busy for more than
a quarter of the run, or if CPU steal was over one percent.
\item A cell is one simulator version on one host, running one case of
one design at one cycle cap. It reports the median of its kept runs, with
their minimum and maximum, of compile time, run time, their total, cycles
per second of the run, and peak memory of the run.
\item A cell that loses more than a third of its runs is not reported; it
needs more runs.
\item A cell with a failing run is reported as failing.
\item A rate from runs under one second, or with a run that the timer
measured as zero, is marked with a dagger ($\dagger$): wall time is
counted in hundredths of a second. With a run measured as zero there is
no rate at all, which the tables show as ``n/a''.
\item A cell whose runs simulated different numbers of cycles is not
summarized.
\end{itemize}
For short simulations the total decides which simulator is faster; for
long ones, cycles per second.

\section{Results}

The tables follow on landscape pages, one per campaign, host and
language. A dagger ($\dagger$) marks a resolution-limited rate.

\begin{landscape}
\input{tables}
\end{landscape}

\section{Provenance}

The cells ran at these commits of the repository, with these lockfiles:

\input{provenance}

The raw per-run files, the command lines of every run, and each
campaign's notes are released with this report.

\end{document}
